\section{Sumcheck As Claim Randomization}

Our target is the MLE evaluation-handle interface.  After the prover is bound
to a tensor \(h\in\F^{\B^d}\), the verifier should be able to check evaluation
claims of the form
\[
  (\widehat O_h,z,v)\longmapsto 1
  \qquad\text{iff}\qquad
  \MLE(h)(z)=v,
\]
where \(\widehat O_h\) is the native MLPCS handle, or the ideal handle being
realized by a concrete PCS.  Sumcheck does not implement this functionality by
itself.  Instead, it transforms one evaluation claim about the same committed
or idealized object into another evaluation claim:
\[
  (\widehat O_h,z,v)\quad\rightsquigarrow\quad(\widehat O_h,r,u).
\]
This does not remove the need to verify the new claim
\((\widehat O_h,r,u)\).  The improvement is in the distribution of the point.
The original point \(z\) may be known before the prover commits; the new point
\(r\) is sampled by the verifier after the prover is bound to the handle and to
the sumcheck messages.  As protocol designers, that randomization is valuable.

The underlying sumcheck protocol is the classical one from LFKN92~[LFKN92].
This section emphasizes the multilinear-evaluation and batching viewpoint used
by HyperPlonk, BaseFold, and Blaze~[CBBZ23,BaseFold23,Blaze24].

The first step is to express the claim $f(z)=v$ as a claim about the Boolean
cube.  A multilinear polynomial is determined by its values on $\B^d$.  For
$b\in\B^d$, define the Boolean interpolation weight
\[
  \chi_b(z)
  =
  \prod_{i=0}^{d-1}\bigl(b_i z_i+(1-b_i)(1-z_i)\bigr).
\]
These weights project the arbitrary point $z$ back onto the Boolean cube:
\[
  f(z)
  =
  \sum_{b\in\B^d} f(b)\chi_b(z).
\]
This identity is not an optimization.  It only rewrites the arbitrary-point
claim as a weighted Boolean-cube sum.

For fixed $z$, it is convenient to put the weights into a polynomial whose
variable is the Boolean-cube point being summed over:
\[
  \chi_z(Y)
  =
  \prod_{i=0}^{d-1}\bigl((1-Y_i)(1-z_i)+Y_i z_i\bigr).
\]
At a Boolean point $b$, this gives $\chi_z(b)=\chi_b(z)$.  Therefore, if
\[
  H(Y)=f(Y)\chi_z(Y),
\]
then the PCS claim is equivalent to
\[
  f(z)=v
  \qquad\Longleftrightarrow\qquad
  \sum_{b\in\B^d}H(b)=v.
\]
Even though $f$ and $\chi_z$ are multilinear, their product $H$ is generally
not multilinear: a variable can have degree two.  This is fine.  Sumcheck does
not require multilinearity.  It requires a known bound on the individual degree
of $H$ in each variable, because that bound controls the degree of each
univariate round polynomial.
The verifier could check this by reading all Boolean values of $f$, but that is
the trivial PCS of Figure~\ref{fig:protocol-trivial-pcs} again.  Sumcheck keeps
the Boolean-sum meaning while replacing the endpoint by a random point.

\paragraph{Evaluation and coefficient views.}
The discussion above used the Boolean values
\[
  b\longmapsto f(b),
  \qquad b\in\B^d,
\]
because that is the cleanest way to see why the interpolation weights
$\chi_b(z)$ recover $f(z)$.  This is an \emph{evaluation view}: the tensor is
the table of values of $f$ on the Boolean cube.

The same multilinear polynomial can also be represented by a coefficient tensor
$y\in\F^{\B^d}$:
\[
  f(X_0,\ldots,X_{d-1})
  =
  \sum_{b\in\B^d} y[b]\prod_{i=0}^{d-1}X_i^{b_i}.
\]
This is a different basis for the same object, not a different polynomial.  For
example, when $d=1$,
\[
  f(X)=y[0]+y[1]X,
  \qquad
  (f(0),f(1))=(y[0],y[0]+y[1]).
\]
Either representation determines the other by a linear change of basis.

Sumcheck is easiest to motivate in the evaluation view because it starts with a
Boolean-cube sum.  The concrete BaseFold construction is easiest to fold in the
coefficient view.  For an $(i+1)$-variate multilinear polynomial, splitting the
coefficient tensor along the last Boolean coordinate writes
\[
  f(X_0,\ldots,X_i)
  =
  f_L(X_0,\ldots,X_{i-1})
  +
  X_i f_R(X_0,\ldots,X_{i-1}),
\]
and a challenge $\theta\in\F$ specializes that variable:
\[
  f(X_0,\ldots,X_{i-1},\theta)
  =
  f_L(X_0,\ldots,X_{i-1})
  +
  \theta f_R(X_0,\ldots,X_{i-1}).
\]
This is the coefficient-side version of a fold.  Later, when BaseFold commits to
a coefficient tensor, sumcheck is still talking about the same polynomial
function $f$; only the committed representation has changed.

\subsection{The Two-Variable Case}

Take $d=2$.  The verifier wants to check
\[
  H(0,0)+H(1,0)+H(0,1)+H(1,1)=v.
\]
The prover first sends a univariate polynomial
\[
  g_1(X)=H(0,X)+H(1,X).
\]
The verifier can check the claimed total without evaluating $H$:
\[
  g_1(0)+g_1(1)=v.
\]
If this passes, the verifier samples $r_1\leftarrow\F$, sends it to the prover,
and changes the claim to
\[
  H(0,r_1)+H(1,r_1)=g_1(r_1).
\]
One Boolean coordinate has disappeared.  The prover then sends
\[
  g_0(X)=H(X,r_1).
\]
The verifier checks
\[
  g_0(0)+g_0(1)=g_1(r_1),
\]
samples $r_0\leftarrow\F$, sends it to the prover, and changes the claim to
\[
  H(r_0,r_1)=g_0(r_0).
\]
Now there is no Boolean sum left.  Since $H(Y)=f(Y)\chi_z(Y)$, the last claim is
\[
  f(r_0,r_1)\chi_z(r_0,r_1)=g_0(r_0).
\]
The verifier can compute $\chi_z(r_0,r_1)$ locally.  It only needs one
remaining value of the committed polynomial $f$ at the random point
$(r_0,r_1)$.

The order matters.  The prover sends each $g_i$ before seeing the corresponding
random challenge $r_i$.  A false polynomial identity can agree at a random
field point only with small probability, bounded by its degree over $|\F|$.

\begin{protocolbox}{Protocol: Sumcheck Randomizes An MLE Evaluation Claim}{fig:protocol-sumcheck-randomizes-pcs}
\noindent Public input: \(\widehat O_h:\F^d\to\F\), \(z\in\F^d\), \(v\in\F\).

\smallskip
\begin{enumerate}[leftmargin=0.24in,label=\textbf{S\arabic*.},itemsep=0.08em,topsep=0.1em,parsep=0pt]
  \item \(\verifier\): \(s_d\gets v\).
  \item For \(i=d-1,d-2,\ldots,0\):
        \begin{enumerate}[leftmargin=0.28in,label=\textbf{S\arabic*.\alph*.},itemsep=0.05em,topsep=0.05em,parsep=0pt]
          \item \(\prover\to\verifier\): \(g_i\in\F[X]\).
          \item \(\verifier\): if \(\deg(g_i)>2\), return reject.
          \item \(\verifier\): if \(g_i(0)+g_i(1)\ne s_{i+1}\), return reject.
          \item \(\verifier\): \(r_i\leftarrow\F\).
          \item \(\verifier\to\prover\): \(r_i\).
          \item \(\verifier\): \(s_i\gets g_i(r_i)\).
        \end{enumerate}
  \item \(\verifier\): \(r\gets(r_0,\ldots,r_{d-1})\).
  \item \(\verifier\): if \(\chi_z(r)=0\), return reject.
  \item \(\verifier\): \(u\gets s_0/\chi_z(r)\).
  \item \(\verifier\): \(u_{\mathrm{open}}\gets
        \mathcal F_{\mathrm{MLPCS}}.\mathsf{Eval}(\widehat O_h,r)\).
  \item \(\verifier\): if \(u_{\mathrm{open}}\ne u\), return reject.
  \item \(\verifier\): return \((1,\widehat O_h,r,u)\).
\end{enumerate}
\end{protocolbox}

\begin{protocolbox}{Subroutine: SumCheck With Terminal Output}{fig:subroutine-sumcheck-terminal}
\noindent Public parameters: dimension \(d\) and individual degree bound
\(\delta\).  In Step C2.a, empty coordinate blocks are omitted.

\smallskip
\noindent\(\underline{\SumCheck_d(\verifier:S\in\F,\ \prover:G:\F^d\to\F)\to(a\in\{0,1\},r\in\F^d,T\in\F)}\)
\begin{enumerate}[leftmargin=0.24in,label=\textbf{C\arabic*.},itemsep=0.08em,topsep=0.1em,parsep=0pt]
  \item \(\verifier\): \(s_d\gets S\).
  \item For \(i=d-1,d-2,\ldots,0\):
        \begin{enumerate}[leftmargin=0.28in,label=\textbf{C\arabic*.\alph*.},itemsep=0.05em,topsep=0.05em,parsep=0pt]
          \item \(\prover\): \(g_i(X)\gets
                \sum_{b\in\B^i}G(b_0,\ldots,b_{i-1},X,r_{i+1},\ldots,r_{d-1})\).
          \item \(\prover\to\verifier\): \(g_i\in\F[X]\).
          \item \(\verifier\): if \(\deg(g_i)>\delta\), return
                \((0,0^d,0)\).
          \item \(\verifier\): if \(g_i(0)+g_i(1)\ne s_{i+1}\), return
                \((0,0^d,0)\).
          \item \(\verifier\): \(r_i\leftarrow\F\).
          \item \(\verifier\to\prover\): \(r_i\).
          \item \(\verifier\): \(s_i\gets g_i(r_i)\).
        \end{enumerate}
  \item \(\verifier\): \(r\gets(r_0,\ldots,r_{d-1})\).
  \item \(\verifier\): \(T\gets s_0\).
  \item \(\verifier\): return \((1,r,T)\).
\end{enumerate}
\end{protocolbox}

Figure~\ref{fig:subroutine-sumcheck-terminal} is the reusable form of the same
interaction.  It checks all round equations and returns the terminal claim
\(G(r)=T\).  The subroutine does not authenticate that terminal value by
itself.  In this paper, every nontrivial use of \(\SumCheck_d\) must immediately
tie \(T\) back to already committed MLE handles.

Concretely, suppose the caller has hatted MLE handles
\[
  J_t:\F^{d_t}\to\F\qquad(t\in\{1,\ldots,m\}),
\]
public maps \(\psi_t:\F^d\to\F^{d_t}\), and a public expression
\[
  \Phi:\F^d\times\F^m\to\F
\]
such that the sumcheck polynomial is meant to satisfy
\[
  G(X)=
  \Phi\bigl(X,J_1(\psi_1(X)),\ldots,J_m(\psi_m(X))\bigr).
\]
After \(\SumCheck_d\) returns \((a,r,T)\), the honest prover sends
\[
  w_t\gets J_t(\psi_t(r))\qquad(t\in\{1,\ldots,m\}).
\]
The verifier records the evaluation requests
\[
  \EvalReqs_{\mathrm{term}}
  =
  \{(J_t,\psi_t(r),w_t):t\in\{1,\ldots,m\}\}
\]
and checks
\[
  T=\Phi(r,w_1,\ldots,w_m).
\]
Only after these endpoint requests are answered by the native
\(\mathcal F_{\mathrm{MLPCS}}\) interface, or authenticated by the compiled
backend that realizes that interface, does the sumcheck terminal claim become
meaningful.

The protocol in Figure~\ref{fig:protocol-sumcheck-randomizes-pcs} is
deliberately circular if the endpoint is the ideal MLPCS functionality: we have
transformed one MLE evaluation into another MLE evaluation.  The reason to keep
the transformation is that the new point \(r\) is verifier-random and
transcript-bound.  BaseFold removes the circularity by implementing this random
endpoint from a commitment to an encoded oracle.  In short, sumcheck randomizes
the polynomial claim; BaseFold folds the committed codeword along the same
challenges.

\subsection{Inner-Product Sumcheck}

The same randomization idea also explains the inner-product checks used later by
Blaze.  Let
\[
  w,x\in\F^{\B^\ell}.
\]
For this subsection, both tensors are Boolean-evaluation tensors.  The verifier
knows $w$ and can compute with all of it.  The prover is bound to $x$ by some
commitment handle $C_x$, but the verifier should not have to read all of $x$.
The claim is
\[
  \sum_{b\in\B^\ell} w[b]x[b]=S.
\]
The objective is to move this full inner product to one claim about the
committed tensor $x$ at a verifier-random point.

Define the polynomial
\[
  H(Y)=\MLE(w)(Y)\MLE(x)(Y).
\]
On Boolean inputs, $H(b)=w[b]x[b]$.  Therefore the inner-product claim is
exactly the Boolean-sum claim
\[
  \sum_{b\in\B^\ell}H(b)=S.
\]
This $H$ is not generally multilinear, because it is a product of two
multilinear polynomials.  Its individual degree is at most two, and that is the
degree bound sumcheck needs.  Therefore the round polynomials are degree two.

Trace one round.  Suppose the active tensors are still indexed by $\B^d$ and
the current claim is their inner product:
\[
  \sum_{a\in\B^d} w[a]x[a]=s_d.
\]
To eliminate the last coordinate, write each address as $(a\mid b)$ with
$a\in\B^{d-1}$ and $b\in\B$.  For each $a\in\B^{d-1}$ define the two line
interpolants
\[
  w_a(X)=(1-X)w[(a\mid0)]+Xw[(a\mid1)],
\]
\[
  x_a(X)=(1-X)x[(a\mid0)]+Xx[(a\mid1)].
\]
The honest round polynomial is
\[
  p(X)=\sum_{a\in\B^{d-1}}w_a(X)x_a(X).
\]
This polynomial has a clear job:
\[
  p(0)+p(1)=s_d.
\]
After checking this equation, the verifier samples $r_{d-1}\leftarrow\F$, sends
it to the prover, and
the new claim becomes
\[
  \sum_{a\in\B^{d-1}} w_a(r_{d-1})x_a(r_{d-1})=s_{d-1},
  \qquad s_{d-1}=p(r_{d-1}).
\]
So one round folds both tensors by the same verifier challenge and replaces the
old inner product by a smaller inner product.

The verifier can perform the fold of $w$ locally, because $w$ is public.  It
cannot locally fold every entry of $x$, because $x$ is committed.  That is the
point of the randomization: after $\ell$ rounds, the verifier has folded $w$
down to the public scalar $\MLE(w)(r)$, and the sumcheck transcript has
produced a random point $r=(r_0,\ldots,r_{\ell-1})$ and a final scalar $s_0$.
Assuming $\MLE(w)(r)\ne0$ in this simplified presentation, the verifier sets
\[
  u=s_0/\MLE(w)(r).
\]
Thus the input inner-product claim has been transformed into the ordinary
committed-evaluation claim
\[
  (C_x,r,u).
\]
Here $(C_x,r,u)$ means that the tensor bound to $C_x$ should satisfy
$\MLE(x)(r)=u$.  This output claim is useful because it asks for one random
evaluation of the committed tensor $x$, not for all of $x$.  BaseFold is the
mechanism that later binds that value $u$ to the committed encoding of $x$ with
only a few openings.
